%% ============================================================
%%  §6b  Relational Semantics and Interpreter Adequacy
%% ============================================================

\section{Relational Semantics and Interpreter Adequacy}
\label{sec:semantics}

The \texttt{Translation/ElabSpec.lean} theorems (\S\ref{sec:rust-proofs})
prove that the \emph{fuel-based interpreter} produces correct outputs.
A natural question is whether the interpreter faithfully captures the
\emph{intended meaning} of LLBC programs — i.e., whether it is adequate
with respect to a language-independent semantic ground truth.
This section answers that question via a relational big-step semantics
(\texttt{Translation/Semantics.lean}) and a soundness proof
(\texttt{Translation/Adequacy.lean}).

\subsection{Relational Big-Step Semantics}

We define a standard inductive big-step operational semantics over the
same runtime types (\texttt{Value}, \texttt{Locals}, \texttt{Signal},
\texttt{FunEnv}) as the interpreter — sharing types avoids any
translation layer in the adequacy proof.

The semantics consists of six mutually inductive relations:

\begin{center}
\small
\begin{tabular}{@{}lll@{}}
\toprule
\textbf{Relation} & \textbf{Judgment} & \textbf{Meaning} \\
\midrule
\texttt{EvalPlace}   & $s \vdash p \Downarrow v$           & Place $p$ reads value $v$ from locals $s$ \\
\texttt{EvalOperand} & $s \vdash o \Downarrow v$           & Operand $o$ evaluates to $v$ \\
\texttt{EvalRValue}  & $s \vdash r \Downarrow v$           & R-value $r$ evaluates to $v$ \\
\texttt{EvalWrite}   & $s \xrightarrow{p := v} s'$         & Writing $v$ to place $p$ yields $s'$ \\
\texttt{EvalStmt}    & $\Gamma \vdash s, \sigma \Downarrow \sigma', s'$ & Statement executes with signal $\sigma$ \\
\texttt{EvalFun}     & $\Gamma \vdash f(\overline{v}) \Downarrow v'$ & Function call terminates with $v'$ \\
\bottomrule
\end{tabular}
\end{center}

\noindent
\textbf{Design choices.}
Loops are handled by three rules — \texttt{loop\_cont} (body produces
\texttt{next}; continue looping), \texttt{loop\_break} (body produces
\texttt{break\_}; exit), and \texttt{loop\_return} (body produces
\texttt{return\_}; propagate) — so the relation stays strictly inductive
(no coinduction).
Panics are modelled by \emph{absence} of a derivation: there is no
$\texttt{EvalStmt}$ rule for division by zero or integer overflow.
This matches the Lean convention that partial functions are undefined
on bad inputs rather than mapped to an error value.

\begin{lstlisting}[caption={Selected \texttt{EvalStmt} rules (Lean notation).},
                   label={lst:evalstmt}]
-- Sequential composition: s1 produces next, then s2 runs
| seq_next {s1 s2 sig s s_mid s'} :
    EvalStmt env s1 s .next s_mid →
    EvalStmt env s2 s_mid sig s' →
    EvalStmt env (.seq s1 s2) s sig s'

-- Loop continues: body produces .next
| loop_cont {body s s' s''} :
    EvalStmt env body s .next s' →
    EvalStmt env (.loop body) s' sig s'' →
    EvalStmt env (.loop body) s sig s''

-- Loop breaks: body produces .break_
| loop_break {body s s'} :
    EvalStmt env body s .break_ s' →
    EvalStmt env (.loop body) s .next s'
\end{lstlisting}

\subsection{Soundness}

The main adequacy theorem asserts that every output produced by the
fuel-based interpreter is derivable in the relational semantics:

\begin{lstlisting}[caption={Soundness of the LLBC interpreter.},
                   label={lst:soundness}]
theorem evalStmtFuel_sound
    {env : FunEnv} {n : Nat} {stmt : LLBCStmt}
    {s : Locals} {sig : Signal} {s' : Locals}
    (h : evalStmtFuel env n stmt s = Except.ok (sig, s')) :
    EvalStmt env stmt s sig s' := ...
\end{lstlisting}

\noindent
The proof is by structural induction on $n$ (the fuel counter), with a
case split on the statement constructor.
Each case reduces the hypothesis \texttt{h} to a concrete
\texttt{Except.ok} equation via \texttt{simp} and \texttt{split},
then constructs the corresponding relational derivation.
The proof comprises 14 theorems (0 sorry), covering all LLBC statement forms:
\texttt{assign}, \texttt{return\_}, \texttt{skip}, \texttt{break\_},
\texttt{continue\_}, \texttt{seq} (both next and abort variants),
\texttt{ite}, \texttt{loop} (all three exit modes), and \texttt{call}.

\paragraph{Key proof patterns.}
Two patterns appear throughout the adequacy proof:

\begin{itemize}[leftmargin=2em]
  \item \textbf{\texttt{split at h}} for nested \texttt{match} reductions.
        Using \texttt{rcases} + \texttt{simp} instead leads to lambda
        alpha-renaming issues that break dependent pattern matching; \texttt{split}
        avoids this by keeping the match scrutinee in the context.

  \item \textbf{\texttt{change} for do-notation} reductions.
        Lean's kernel reduces \texttt{do}-notation to nested \texttt{match} on
        \texttt{Except.ok} / \texttt{Except.error}; the \texttt{change}
        tactic asserts the reduced form explicitly, avoiding opaque unfolding
        steps that would otherwise require \texttt{norm\_num} or \texttt{ring}.
\end{itemize}

\subsection{Completeness}

The converse direction — every relational derivation is witnessed by
some fuel bound — is proved as \texttt{evalStmtFuel\_complete}:

\begin{lstlisting}[caption={Completeness of the LLBC interpreter.},
                   label={lst:complete}]
theorem evalStmtFuel_complete
    {env : FunEnv} {stmt : LLBCStmt}
    {s : Locals} {sig : Signal} {s' : Locals}
    (h : EvalStmt env stmt s sig s') :
    ∃ n, evalStmtFuel env n stmt s = Except.ok (sig, s') := ...
\end{lstlisting}

\noindent
The proof proceeds by structural induction on the \texttt{EvalStmt} derivation.
Two auxiliary lemmas are key:

\begin{itemize}[leftmargin=2em]
  \item \textbf{\texttt{evalStmtFuel\_mono\_one}}: if
        \texttt{evalStmtFuel env n stmt s = .ok r}, then
        \texttt{evalStmtFuel env (n+1) stmt s = .ok r}.
        Proved by induction on $n$, case analysis on \texttt{stmt}.

  \item \textbf{\texttt{evalStmtFuel\_mono}}: monotonicity for all $k$,
        derived from \texttt{evalStmtFuel\_mono\_one} by induction on $k$.
\end{itemize}

\noindent
For the \texttt{seq\_next} case, the existential witnesses are
$\max(n_1, n_2) + 1$ where $n_1, n_2$ come from the two inductive
hypotheses; \texttt{evalStmtFuel\_mono} is used to lift each sub-proof
to the common bound.
The \texttt{loop\_cont} case requires composing the monotonicity lemma for
the body and the recursive loop call; this is the most technically intricate
case but follows the same pattern.

Together, \texttt{evalStmtFuel\_sound} and \texttt{evalStmtFuel\_complete}
establish the \emph{full adequacy} of the interpreter: the fuel-based
operational semantics and the inductive big-step semantics are extensionally
equivalent. The \texttt{Translation/Adequacy.lean} file contains
\textbf{19 theorems (0 \sorry)}.

\subsection{Relationship to the Verification Pipeline}

The relational semantics serves as the \emph{semantic ground truth} against
which the interpreter is calibrated:

\[
  \underbrace{\texttt{evalStmtFuel}}_{\text{interpreter}}
  \;\overset{\text{sound}}{\longrightarrow}\;
  \underbrace{\texttt{EvalStmt}}_{\text{relational semantics}}
\]

The \texttt{ElabSpec} and \texttt{CharonSpec} theorems operate at the
interpreter level. Soundness guarantees that any property proved there
also holds in the relational model. This closes the semantic gap between
the constructive evaluation approach (easy to compute with, easy to prove
concrete instances by \texttt{rfl}) and the declarative inductive semantics
(easy to reason about symbolically, standard in PL theory).
